\rtmtight
\begin{tblr}{width=\textwidth,colspec={|Q[c,m,2.1cm]|X[l,m]|},hlines,vlines,hspan=minimal,rowsep=1.5pt,colsep=3pt,row{1}={bg=headgray,font=\bfseries}}
\SetCell{c,m}\hfil\logo\hfil & \textbf{POLITEKNIK NEGERI MALANG}\par\textbf{JURUSAN TEKNOLOGI INFORMASI}\par\textbf{PROGRAM STUDI: D4 TEKNIK INFORMATIKA} \\
\SetCell[c=2]{l} \textbf{RENCANA TUGAS MAHASISWA (RTM) DAN RUBRIK PENILAIAN} & \\
Nama Mata Kuliah & Basis Data \\
Kode Mata Kuliah & RTI252006 \quad \textbf{sks / Jam perminggu:} 2 SKS/4 jam \quad \textbf{Semester:} 2 \\
Dosen Pengampu & \begin{enumerate}\item Candra Bella Vista, S.Kom., M.T.\item Elok Nur Hamdana, S.T., M.T.\item M. Hasyim Ratsanjani, S.Kom., M.Kom.\item Ridwan Rismanto, S.ST., M.Kom., Ph.D.\item Vit Zuraida, S.Kom., M.Kom.\end{enumerate} \\
\end{tblr}
\assessmentblock{Tugas Perancangan dan Implementasi Basis Data}{Tugas terstruktur berbasis kasus (CBL)}{SCPMK0203-01501, SCPMK1001-01501}{Mahasiswa menyelesaikan rangkaian Tugas 1--9: mengulas penerapan basis data, menganalisis kebutuhan data, merancang ERD, memetakan ke model relasional, melakukan normalisasi, serta menuliskan perintah SQL DDL/DML/DQL sesuai studi kasus.}{Menganalisis kebutuhan data studi kasus, merancang ERD dan model relasional, melakukan normalisasi, mengimplementasikan rancangan pada MySQL, dan mendokumentasikan hasil setiap tugas.}{Lembar jawaban tugas, diagram ERD, model relasional, skrip SQL, dan dokumentasi hasil eksekusi.}{Ketepatan analisis kebutuhan data dan rancangan ERD & 5 & Entitas, atribut, relationship, kardinalitas, dan partisipan diidentifikasi lengkap dari studi kasus dan ERD konsisten dengan kebutuhan data. & Komponen utama ERD benar, tetapi sebagian kardinalitas, partisipan, atau atribut kurang tepat. & ERD tidak sesuai kebutuhan data studi kasus atau banyak komponen salah/tidak lengkap. \\ \\
Ketepatan pemetaan model relasional dan normalisasi & 5 & Pemetaan ERD ke model relasional mengikuti algoritma dengan benar dan hasil normalisasi mencapai 3NF tanpa anomali. & Pemetaan dan normalisasi sebagian besar benar, tetapi ada kesalahan pada relasi tertentu atau normalisasi berhenti di 2NF. & Pemetaan tidak mengikuti algoritma atau hasil normalisasi masih menyisakan redundansi dan anomali. \\ \\
Ketepatan perintah SQL-DDL & 5 & Perintah CREATE, ALTER, dan DROP lengkap, rapi, dan berhasil mengimplementasikan rancangan termasuk primary key dan foreign key. & Perintah DDL sebagian besar benar, tetapi ada kesalahan pada tipe data, constraint, atau relasi. & Perintah DDL tidak dapat dieksekusi atau tidak sesuai rancangan basis data. \\ \\
Ketepatan perintah SQL-DML dan DQL & 5 & Perintah INSERT, UPDATE, DELETE, dan SELECT (termasuk klausa dan fungsi agregasi) benar dan sesuai kebutuhan studi kasus. & Perintah DML/DQL sebagian besar benar, tetapi ada klausa atau kondisi yang kurang tepat. & Perintah DML/DQL salah, tidak dapat dieksekusi, atau tidak menjawab kebutuhan studi kasus. \\}

\assessmentblock{Kuis Konsep dan SQL Basis Data}{Kuis (ujian teori pilihan ganda dan/atau esai, closed book)}{SCPMK0203-01501, SCPMK1001-01501}{Mahasiswa mengerjakan Kuis 1 (materi minggu 1--3: konsep basis data dan perancangan ERD) dan Kuis 2 (materi minggu 10--13: SQL DDL, DML, DQL, dan JOIN) secara mandiri dan jujur.}{Mengerjakan soal pilihan ganda dan/atau esai secara individu dengan sifat closed book sesuai jadwal asesmen.}{Lembar jawaban kuis.}{Kuis 1: penguasaan konsep basis data dan perancangan ERD & 10 & Jawaban tepat sesuai kunci; konsep basis data relasional, pemodelan data, dan komponen ERD dijelaskan benar pada kasus. & Sebagian besar jawaban benar, tetapi ada kekeliruan pada penerapan konsep ke kasus. & Banyak jawaban tidak sesuai kunci atau konsep dasar basis data belum dikuasai. \\ \\
Kuis 2: penguasaan SQL DDL, DML, dan DQL & 10 & Jawaban tepat sesuai kunci; sintaks dan penggunaan perintah SQL (CREATE, INSERT, UPDATE, DELETE, SELECT, JOIN) benar sesuai kasus. & Sebagian besar jawaban benar, tetapi ada kesalahan sintaks atau pemilihan perintah. & Banyak jawaban salah atau perintah SQL yang dituliskan tidak sesuai kebutuhan kasus. \\}

\assessmentblock{UTS Perancangan Basis Data}{UTS (ujian teori pilihan ganda dan/atau esai, closed book)}{SCPMK0203-01501}{Mahasiswa mengerjakan soal UTS materi minggu 1--8 meliputi konsep basis data relasional, perancangan ERD, pemetaan ke model relasional, dan normalisasi secara mandiri dan jujur.}{Mengerjakan soal pilihan ganda dan/atau esai secara individu dengan sifat closed book sesuai jadwal asesmen.}{Lembar jawaban UTS.}{Penguasaan konsep data dan basis data relasional & 8 & Jawaban tepat sesuai kunci; pengertian, karakteristik, jenis, dan prinsip pengembangan basis data relasional dijelaskan benar. & Sebagian jawaban benar, tetapi penjelasan konsep pada kasus belum lengkap. & Banyak jawaban tidak sesuai kunci atau konsep dasar tidak dikuasai. \\ \\
Ketepatan perancangan ERD dan pemetaan ke model relasional & 12 & Rancangan ERD dan hasil pemetaan ke model relasional benar sesuai algoritma, termasuk kardinalitas dan partisipan. & Rancangan dan pemetaan sebagian besar benar, tetapi ada kesalahan pada relasi atau kardinalitas tertentu. & Rancangan ERD atau pemetaan tidak sesuai kebutuhan data dan algoritma pemetaan. \\ \\
Ketepatan proses dan hasil normalisasi & 10 & Tahapan 1NF--3NF diterapkan benar pada tabel soal dan hasil akhir bebas redundansi dan anomali. & Tahapan normalisasi sebagian benar, tetapi hasil akhir belum sepenuhnya 3NF. & Tahapan normalisasi salah atau hasil masih menyisakan banyak anomali. \\}

\assessmentblock{Case Method Studi Kasus Basis Data}{Case Method}{SCPMK1001-01502}{Mahasiswa menyelesaikan studi kasus minggu 15--16: mengidentifikasi kebutuhan data organisasi/bisnis, merancang ERD, memetakan ke model relasional (3NF), mengimplementasikan ke MySQL, dan menyusun query DML/SELECT.}{Menganalisis kebutuhan data studi kasus secara berkelompok, merancang dan menormalisasi skema, mengimplementasikan pada DBMS MySQL, menguji query, dan menyusun laporan perancangan dan implementasi.}{Laporan perancangan dan implementasi basis data, file ERD dan model relasional, skrip SQL, serta bukti eksekusi query.}{Ketepatan identifikasi kebutuhan dan rancangan basis data & 4 & Kebutuhan data lengkap (entitas, atribut, relationship, kardinalitas, partisipan) dan ERD serta model relasional 3NF konsisten dengan studi kasus. & Kebutuhan utama teridentifikasi, tetapi sebagian rancangan atau bentuk normal belum tepat. & Kebutuhan data tidak teridentifikasi dengan benar atau rancangan tidak sesuai studi kasus. \\ \\
Ketepatan implementasi basis data pada MySQL & 3 & Query DDL mengimplementasikan seluruh hasil rancangan dengan benar termasuk relasi antartabel, dan DML mengisi data secara konsisten. & Implementasi sebagian besar berhasil, tetapi ada tabel, relasi, atau data yang belum sesuai rancangan. & Implementasi gagal dieksekusi atau tidak sesuai hasil rancangan. \\ \\
Ketepatan query penyajian data dan kualitas laporan & 3 & Query SELECT termasuk multi table menjawab kebutuhan informasi studi kasus dan laporan mendokumentasikan proses secara runtut. & Query utama benar, tetapi sebagian kebutuhan informasi belum terjawab atau laporan kurang lengkap. & Query tidak menjawab kebutuhan informasi atau laporan tidak menggambarkan proses pengerjaan. \\}

\assessmentblock{UAS SQL Studi Kasus}{UAS (ujian teori pilihan ganda dan/atau esai, closed book)}{SCPMK1001-01502}{Mahasiswa mengerjakan soal UAS materi minggu 1--16 dengan penekanan pada penyusunan SQL yang berkaitan dengan studi kasus secara mandiri dan jujur.}{Mengerjakan soal pilihan ganda dan/atau esai secara individu dengan sifat closed book sesuai jadwal asesmen.}{Lembar jawaban UAS.}{Ketepatan analisis kasus dan rancangan skema & 6 & Analisis kebutuhan data soal kasus tepat dan skema/model relasional yang disusun konsisten dengan kasus. & Analisis dan skema sebagian besar benar, tetapi ada kebutuhan data yang terlewat. & Analisis kasus salah atau skema tidak sesuai kebutuhan data. \\ \\
Ketepatan penyusunan query SQL sesuai studi kasus & 10 & Query DDL, DML, dan SELECT (termasuk JOIN, klausa, dan agregasi) benar sintaks dan tepat menjawab kebutuhan kasus. & Sebagian besar query benar, tetapi ada kesalahan sintaks atau kondisi pada beberapa soal. & Query banyak yang salah sintaks atau tidak menjawab kebutuhan kasus. \\ \\
Efisiensi dan kerapian penulisan query & 4 & Query ditulis efisien, rapi, dan menggunakan klausa yang sesuai tanpa langkah berlebihan. & Query berfungsi tetapi kurang efisien atau penulisannya kurang rapi. & Query berbelit, tidak efisien, dan sulit dibaca. \\}

\begin{tblr}{width=\textwidth,colspec={|Q[l,m,3.2cm]|X[l,m]|},hlines,vlines,hspan=minimal,row{1-3}={bg=headgray},rowsep=1.5pt,colsep=3pt}
Jadwal Pelaksanaan: & Minggu 1--17 sesuai RPS. Setiap evaluasi memiliki RTM dan rubrik multi-kriteria. \\
Lain-Lain Yang Diperlukan: & LMS, DBMS MySQL/MariaDB, perangkat pemodelan ERD, dan perangkat presentasi. \\
Pustaka: & Mengacu pustaka utama dan pendukung pada bagian RPS. \\
\end{tblr}
